%% kapa: exam=1 round=26 year=2015 period=2 form=A subject=appraisal_law subject_ko=부동산관계법규 seq=8 number=8
%% source: Q-Net 감정평가사 기출문제 (cst00309) artlSeq=5138001 file=extracted/1008/5138001/제26회_감정평가사 1차_기출/제26회_감정평가사 1차_기출/2교시_감평_A형.pdf
%% conversion: kordoc 4.24.51 -> tools/gen1.py (auto); review=proofread-sonnet; flags=none
\begin{question}{8}
국토의 계획 및 이용에 관한 법령상 용도지구에 관한 설명으로 옳지 않은 것은?
\begin{choices}
  \item 용도지구란 토지의 이용 및 건축물의 용도ㆍ건폐율ㆍ용적률ㆍ높이 등에 대한 용도지역의 제한을 강화하거나 완화하여 적용함으로써 용도지역의 기능을 증진시키고 미관ㆍ경관ㆍ안전 등을 도모하기 위하여 도시ㆍ군관리계획으로 결정하는 지역을 말한다.
  \item 국토교통부장관은 지정된 용도지구의 일부에 대하여 지구단위계획구역을 지정할 수 없다.
  \item 국토교통부장관은 미관지구를 중심지미관지구, 역사문화미관지구, 일반미관지구로 세분하여 지정할 수 있다.
  \item 건축물이나 그 밖의 시설의 용도ㆍ종류 및 규모 등의 제한은 해당 용도지구의 지정목적에 적합하여야 한다.
  \item 지방자치단체의 장이 다른 법률에 따라 토지 이용에 관한 지구를 지정하는 경우에는 그 지구의 지정목적이「국토의 계획 및 이용에 관한 법률」에 따른 용도지구의 지정목적에 부합되도록 하여야 한다.
\end{choices}
\end{question}
\answer{2}
